\section{Related Work}

\paragraph{Refactoring and software evolution.}
Classical refactoring work established both the breadth of refactoring activities and their importance for long-term software maintenance~\cite{Mens2004RefactoringSurvey}. Later empirical studies showed that refactoring decisions arise regularly in open-source projects and span a wide range of structural transformations rather than a handful of stylized edits~\cite{Kim2012RefactoringInTheWild,Tsantalis2014RefactoringOperations}. More generally, software-engineering research has emphasized that software artifacts co-evolve over time, and that understanding maintenance behavior often requires repository- or revision-oriented analysis rather than one-shot task evaluation~\cite{Zaidman2010CoEvolutionProductionTest}. P1 builds directly on this maintenance lens, but shifts the target artifact from code-and-tests alone to verified programs whose code, contracts, and proof hints evolve together.

\paragraph{Verification engineering and proof maintenance.}
Work on proof maintenance has mostly appeared in theorem-proving settings. Whiteside et al.\ formalize proof-script refactoring as a semantics-preserving restructuring activity~\cite{Whiteside2011ProofScriptRefactoring}, and Whiteside's later thesis expands this line into a larger framework for refactoring proofs~\cite{Whiteside2013RefactoringProofs}. Tool-oriented work such as Tactician shows that proof refactoring can also be operationalized in interactive settings~\cite{Adams2015Tactician}. In the Dafny ecosystem, adjacent work has instead focused on proof guidance and verification support: reusable verification patterns~\cite{Grov2016VerificationPatternsDafny}, testing and differential validation of the Dafny toolchain~\cite{Irfan2022TestingDafny,Fedchin2023ToolkitTestingDafny}, and empirical studies of how auto-active verification affects development practice~\cite{Mugnier2025ImpactVerification}. These papers make two points clear: proof-oriented artifacts are structured enough to merit systematic support, and maintainability matters in practice. However, they do not provide a benchmark for refactoring-induced proof brittleness in an SMT-backed verifier like Dafny.

\paragraph{Automated repair and learning-based assistance.}
Automated software repair supplies the closest high-level analogy for this paper's baseline design, especially generate-and-validate workflows and search over candidate fixes~\cite{Gazzola2017AutomaticSoftwareRepair,Xia2023APRWithPLM}. On the proof side, systems such as TacTok and Passport show that proof corpora can support meaningful automation and synthesis~\cite{First2020TacTok,SanchezStern2023Passport}. Meanwhile, pretrained and generative code models such as CodeBERT, GraphCodeBERT, and InCoder have become standard reference points for learning-based code generation and repair~\cite{Feng2020CodeBERT,Guo2020GraphCodeBERT,Fried2022InCoder}. Recent Dafny-specific work extends this trend toward verified code generation and benchmarking, including LLM-assisted synthesis of verified Dafny methods~\cite{Misu2024AIAssistedDafny}, DafnyBench~\cite{DafnyBench2024}, DafnyComp~\cite{Xu2026DafnyComp}, and Re:Form's RL-oriented formal verification pipeline~\cite{Yan2025ReFormDafny}. This paper uses that literature primarily as baseline and context: it does not propose a repair engine, but rather quantifies how much of the maintenance problem remains after simple baselines such as re-verification, syntactic transfer, and LLM-only repair.

\paragraph{Gap and positioning of P1.}
Dafny has long been used in benchmark-style verification studies~\cite{Leino2010Dafny,Leino2010VSTTEBenchmarks}, and recent work shows growing interest in Dafny-oriented benchmarks for synthesis, compositional reasoning, testing, and benchmark generation~\cite{DafnyBench2024,Wang2026TACoDafny,Xu2026DafnyComp}. Likewise, prior work has studied refactoring broadly, proof refactoring in interactive provers, and repair or synthesis through both symbolic and learned techniques. What is still missing is a \emph{maintenance-centred} benchmark and empirical study centered on \emph{semantics-preserving refactoring of already verified Dafny programs}, together with baseline evidence on how much nearby-version proof breakage simple recovery strategies can repair. This paper targets that missing slice of the space and, by doing so, provides a reusable evaluation substrate for later repair work without collapsing benchmark construction into repair-system design.
